\chapter{Marasmus}
\index{Marasmus}
\label{sec:Marasmus}

\section{Overview}
Marasmus is a form of severe protein-energy malnutrition characterized by inadequate intake of both protein and calories. It is most common in infants and young children in low-resource settings, particularly during the first year of life, and is linked to early weaning, repeated infections, and poverty.
\section{Causes}
This condition happens because of catabolism of muscle and subcutaneous fat to meet energy needs, leading to profound wasting with relative preservation of serum albumin.     
\section{Risk Factors}

\begin{itemize}[noitemsep]
  \item Genetic predisposition and family history
  \item Advancing age
  \item Lifestyle factors (diet, exercise, smoking, alcohol)
  \item Environmental exposures
  \item Underlying medical conditions
  \item Medication use
\end{itemize}
\section{Signs and Symptoms}

\subsection{Common symptoms}
\begin{itemize}[noitemsep]
  \item You may experience severe wasting of muscle and subcutaneous tissue, a ``monkey-like'' facies with prominent eyes, wrinkled skin, irritability, and growth retardation
  \item Subcutaneous fat is absent and there is no peripheral swelling.
  \item Pain or discomfort in the affected area
  \end{itemize}

\subsection{Emergency warning signs}
\begin{itemize}[noitemsep]
  \item Severe or worsening symptoms of marasmus require immediate medical evaluation
\end{itemize}
\section{Progression and Stages}
The condition may progress through different stages of severity over time. Early stages often involve mild or subtle symptoms that may be overlooked. Moderate stages involve more noticeable symptoms affecting daily function. Advanced stages may involve significant impairment and complications.
\section{Complications}
\begin{itemize}[noitemsep]
  \item Disease progression leading to worsening symptoms
  \item Secondary infections or injuries
  \item Side effects of treatment
  \item Reduced quality of life
  \item Emotional and psychological impact
\end{itemize}
\section{Diagnosis}
\begin{itemize}[noitemsep]
  \item Doctors diagnose this based on anthropometric criteria: weight-for-age less than 60\% of median, weight-for-height less than 70\% of median, or mid-upper arm circumference (MUAC) less than 115 mm in children aged 6--59 months. Management follows the WHO protocol: stabilization with correction of dehydration and electrolyte imbalances (using ReSoMal), empiric antibiotics, micronutrient supplementation, and staged refeeding with F-75 formula followed by F-100 formula
  \item Slow, cautious refeeding is required to avoid refeeding syndrome.
  \item Comprehensive medical history and physical examination
  \item Laboratory tests to assess organ function
  \item Imaging studies to visualize affected structures
  \item Specialized testing depending on the affected system
  \item Consultation with relevant specialists
\end{itemize}
\section{Prevention}
\begin{itemize}[noitemsep]
  \item Maintain a healthy lifestyle with balanced diet and exercise
  \item Avoid known risk factors
  \item Attend regular medical check-ups for early detection
  \item Follow recommended screening guidelines
\end{itemize}
\section{Treatment Options}
Medications to manage symptoms and address underlying mechanisms Lifestyle modifications and self-care measures Physical therapy and rehabilitation Surgical intervention when indicated Regular monitoring and follow-up care Patient education and support services

\begin{itemize}[noitemsep]
  \item Medications to manage symptoms and address underlying mechanisms
  \item Lifestyle modifications and self-care measures
  \item Physical therapy and rehabilitation
  \item Surgical intervention when indicated
  \item Regular monitoring and follow-up care
\end{itemize}
\section{Living with It}
\begin{itemize}[noitemsep]
  \item Follow your treatment plan as prescribed
  \item Maintain regular follow-up with your healthcare provider
  \item Make necessary lifestyle adjustments
  \item Seek support from family, friends, or support groups
  \item Stay informed about your condition
\end{itemize}
\section{Outlook}
The outlook is good with timely treatment, but mortality remains high (20--30\%) in severe cases without adequate care. The outlook varies depending on the specific condition, severity at diagnosis, and response to treatment. Early intervention generally leads to better outcomes. Many conditions can be effectively managed with appropriate treatment. Regular follow-up care is important for monitoring and treatment adjustment.
